\begin{tikzpicture}[scale=1]

% === Parâmetros ===
\def\a{2}   % limite das retas x = \pm a

  \begin{axis}[      
      axis lines=middle,            % Eixos no meio da figura
      xlabel={$x$}, ylabel={$z$},
      axis line style={thick},      % Eixos mais grossos
      axis equal image,             % mantém escala x=z para os 45° ficarem corretos
      set layers=axis on top,       % Eixos na camada de cima, para não ser encoberto pelo sombreado
      enlargelimits,                % Pequeno aumento nos limites para melhorar a estética
      xmin=-1.3*\a, xmax=1.3*\a, ymin=-1.3*\a, ymax=1.3*\a,
      clip=false,                   % permite que rótulos ultrapassem a área dos eixos
      xtick={-\a,\a}, ytick={-\a,\a},
      x tick label style={anchor=north east}, % rótulos abaixo e a esquerda do eixo X
      y tick label style={anchor=south east}, % rótulos acima e a esquerda do eixo y
      % width=8cm, height=6cm,
      every axis x label/.style={at={(ticklabel* cs:0.97)}, anchor=north},
      every axis y label/.style={at={(ticklabel* cs:0.97)}, anchor=east},
    ]

    % --- Linhas Tracejadas ---
     \addplot[Tracejado] coordinates {(\a,0) (\a,\a)};
     \addplot[Tracejado] coordinates {(0,\a) (\a,\a)};

    % --- Curvas ---

    \addplot[thick, blue, domain=-1.2*\a:1.2*\a] ({\a}, {x}) node[pos=0.85, right, blue] {$x=2$};   % Reta vertical x=2
    \addplot[thick, blue, domain=-1.2*\a:1.2*\a] ({-\a}, {x}) node[pos=0.85, left, blue] {$x=-2$};  % Reta vertical x=-2

    \addplot[thick, red, domain=-1.2*\a:1.2*\a, samples=100] {x} node[pos=1, above, red] {$z=x$};   % Reta z=x
    \addplot[thick, red, domain=-1.2*\a:1.2*\a, samples=100] {-x} node[pos=1, below, red] {$z=-x$}; % Reta z=-x

    % --- Região sombreada: triângulo no 1o quadrante ---
    \path[fill=gray!30, opacity=0.6]
      (axis cs:0,0) -- (axis cs:\a,0) -- (axis cs:\a,\a) -- cycle;

    % --- Pontos de interseção ---
    \addplot[only marks, black!60, mark=*] coordinates {(\a,\a) (-\a,\a) (\a,-\a) (-\a,-\a)};


  \end{axis}
  
\end{tikzpicture}
